% Philippica XI (philippic-11) --- English translation
% Translated by Claude (Anthropic) from the Latin (Perseus).
% Style guide: STYLE.md.

\ciceroSection{1} In the great grief, senators --- or rather the mourning --- which we have received from the cruel and pitiable death of Gaius Trebonius, that most excellent citizen and most temperate of men, there is nonetheless something which I think will profit the commonwealth. For we have seen plainly how great a savagery dwells in those who have taken up criminal arms against their country. For these two are the foulest and most loathsome monsters born since the human race began: Dolabella and Antony --- of whom the one has accomplished what he longed for, while of the other it has been laid bare what he was plotting. Lucius Cinna was cruel, Gaius Marius unrelenting in his wrath, Lucius Sulla violent; yet the bitterness of none of these in taking vengeance went beyond death --- a penalty which even so was thought too cruel to inflict on citizens.

\ciceroSection{2} Behold, then, a twin pair in crime --- unexampled, unheard of, savage, barbarous. And so those between whom you remember there was once the bitterest hatred and warfare, the very same men were afterwards bound together in a singular accord and affection by the likeness of their utterly depraved nature and their most shameful life. Therefore what Dolabella has done to the one he had in his power, the same Antony threatens to many. But Dolabella --- since he was far from our consuls and our armies, and had not yet perceived that the Senate had joined in concord with the Roman people --- relying on the forces of Antony, undertook those crimes which he supposed had already been undertaken at Rome by the partner of his madness.

\ciceroSection{3} What else, then, do you suppose this man is contriving, what is he longing for --- or what cause for war is there at all? All of us who have held free opinions about the commonwealth, who have spoken sentiments worthy of ourselves, who have wished the Roman people to be free, he has reckoned not as personal foes but as public enemies; yet he plots for us punishments greater than for an enemy: he thinks death the penalty of nature, but torture and torment the penalty of his wrath. What sort of enemy, then, must he be held to be, in whose victory --- if torture were absent --- death would be counted among his kindnesses? Therefore, senators, although you have no need of one to urge you on --- for of your own accord you have caught fire with the longing to recover liberty --- defend that liberty with all the greater spirit and zeal, the greater the punishments of slavery you see set before the conquered.

\ciceroSection{4} Antony has invaded Gaul, Dolabella Asia --- each a province not his own. Against the one Brutus threw himself, and at the peril of his own life he checked the onset of that raving creature, who longed to harry and plunder everything; he barred his advance, reined him back from return, and --- allowing himself to be besieged --- hemmed Antony in on both sides. The other burst into Asia. Why? If it was to reach Syria, the road lay open, sure, and not long; but if it was to reach Trebonius, what need was there of a legion? Having sent ahead one Octavius, some Marsian or other, a criminal brigand and a beggar, to lay waste the fields and harry the cities --- not in any hope of building up an estate, which those who know him deny he is capable of holding (for to me this senator is unknown), but for the immediate feeding of his own destitution --- Dolabella followed after.

\ciceroSection{5} With no suspicion of war --- for who would have thought it? --- there followed the most intimate conversations with Trebonius; embraces of the deepest goodwill stood forth as false witnesses in a feigned affection; right hands that used to be the pledges of good faith were profaned by treachery and crime: an entry by night into Smyrna, as though into an enemy city --- a city of the most faithful and most ancient of allies; Trebonius was overpowered --- if as by one who was openly an enemy, then off his guard; if as by one who still wore even then the appearance of a fellow citizen, then a wretched victim. By this, surely, Fortune wished us to take a lesson of what the conquered must dread. A man of consular rank, holding the province of Asia with consular command, he handed over to a Samian exile: he was unwilling to kill him at once when captured --- lest, I suppose, he should seem too generous in victory. When he had mangled that excellent man with insults of words from his unclean mouth, then by stripes and tortures he conducted his inquiry after the public money, and that for two days together. Afterwards, breaking his neck, he cut off his head, and ordered it to be carried about, fixed upon a pike; the rest of the body, dragged and torn, he flung into the sea.

\ciceroSection{6} With this enemy we must wage war --- an enemy whose most loathsome cruelty has surpassed all barbarism. Why should I speak of the slaughter of Roman citizens, of the plundering of shrines? Who is there who could lament such calamities in a measure worthy of their atrocity? And now he roams over all Asia, flits about like a king; he supposes that we are held off by another war --- as if it were not in fact one and the same war against this accursed yoke of impious men. You behold in Dolabella the very image of Marcus Antonius's cruelty: from him it was moulded; from him the lessons of Dolabella's crimes were handed down. Do you suppose that Antony would be any gentler in Italy, if he were given the chance, than Dolabella was in Asia? For my part, I think both that the one has gone as far as the madness of a savage man could advance, and that Antony, should he gain the power, would leave out no part of any punishment that could be inflicted.

\ciceroSection{7} Set before your eyes, then, senators, that sight --- pitiable indeed and lamentable, but necessary for kindling our spirits: the attack by night upon the most renowned city of Asia, the bursting in of armed men into the house of Trebonius, when the poor man saw the swords of the brigands before he heard what the matter was; the entrance of the raving Dolabella, that foul voice and that infamous face, the chains, the stripes, the rack, the Samian torturer and executioner. They say that he bore these things bravely and with endurance. Great praise --- and in my judgment the greatest of all. For it is the part of a wise man to think out beforehand that whatever may befall a human being is to be borne with composure, should it come to pass. To foresee that no such thing shall happen takes, to be sure, the greater wisdom; to bear it bravely takes no smaller spirit.

\ciceroSection{8} And Dolabella, indeed, was so forgetful of humanity --- though he was never any partaker of it --- that he worked his insatiable cruelty not only upon the living man but even upon the dead; and in mangling and abusing his body, when he could not glut his spirit, he fed his eyes upon it. O Dolabella, far more wretched than the man you wished to be most wretched of all! ``Trebonius endured great pains.'' Many men endure greater from the severity of disease --- men whom, even so, we are accustomed to call not wretched but sorely tried. ``The pain was long.'' Two days' worth --- whereas for many it is often the pain of many years together. Nor in truth are the torments inflicted by executioners worse than are, at times, the torments of disease.

\ciceroSection{9} There are other things --- other things, I say, you most abandoned and most deranged of men --- far more wretched. For inasmuch as the force of the mind is greater than that of the body, by so much are the things conceived in the mind more grievous than those endured in the body. More wretched, then, is he who takes a crime upon himself than he who is compelled to undergo the deed of another. Trebonius was tortured by Dolabella; and Regulus, too, by the Carthaginians. Therefore, since the Carthaginians were judged most cruel for their treatment of an enemy, what is to be judged of Dolabella in his treatment of a fellow citizen? Or is this really a thing to be compared, or to be left in doubt --- which of the two is the more wretched: he whose death the Senate and the Roman people long to avenge, or he who has been judged a public enemy by every vote of the Senate? For in all the other passages of life, who is there who could compare the life of Trebonius with that of Dolabella without the gravest insult to Trebonius? The wisdom, the genius, the humanity, the integrity, the greatness of soul of the one in liberating his country --- who does not know it? To the other, from boyhood, cruelty was a delight; and then so foul a depravity of lusts that he himself always took joy in this very thing: that he did the deeds which could not be cast in his teeth even by a modest enemy.

\ciceroSection{10} And this man --- immortal gods! --- was once mine. For his vices were hidden from one who did not search them out. Nor perhaps would I be estranged from him even now, had he not been found an enemy to us, to the walls of our country, to this city, to the household gods, to the altars and hearths of us all --- in short, to nature and to humanity itself. Warned by his example, let us guard the more carefully and the more watchfully against Antony. For Dolabella did not keep with him so very many notorious and conspicuous brigands; but you see what men, and how many, Antony has. First, his brother Lucius: what a firebrand, immortal gods, what an outrage, what a crime, what a gulf, what an abyss! What do you suppose he would not swallow up in his greed, what not drain off in his designs, whose blood would he not drink, upon whose possessions and fortunes would he not fasten his most shameless eyes in hope and in intent?

\ciceroSection{11} What of Censorinus? --- who said in words that he wished to be urban praetor, but in fact assuredly did not. What of Bestia? --- who openly declares that he is standing for the consulship, in Brutus's place. And this detestable omen may Jupiter avert! But how absurd that a man who could not be made praetor should stand for the consulship --- unless perhaps he counts his conviction in place of a praetorship. The other, the famous Caesar Vopiscus, a man of the highest genius and the highest power, who seeks the consulship straight from the aedileship --- let him be released from the laws: although the laws do not hold him, on account, I suppose, of his extraordinary worth. But this man, with me defending him, was acquitted five times over: a sixth prize in the city is hard even for a gladiator. Yet this is the fault of the jurors, not mine. I defended him in the best of faith: it was for them to keep so distinguished and excellent a senator in the state. And yet he now seems to be doing nothing else than to make us understand that those whose verdicts we annulled judged rightly and for the good of the commonwealth.

\ciceroSection{12} Nor is this true of him alone: there are others in the same camp, honourably condemned and shamefully restored. What plan of these men --- who are enemies to all good men --- do you suppose will prove anything but most cruel? There is added a certain Saxa, whom Caesar gave us as tribune of the plebs from the farthest reaches of Celtiberia, once a marker-out of camps, now, as he hopes, of the city: but since he is a stranger to it, let him pronounce that omen over his own head, and leave us unharmed. With him goes the veteran Cafo, than whom the veterans hate no man worse. To these men Antony --- as though over and above the dowry they had received amid the evils of civil war --- made a lavish grant of the Campanian land, that they might have it as a nurse for the rest of their estates. Would that they were content with it! We should bear it --- though it was not to be borne; but anything had to be suffered, so that we might not have this most loathsome war.

\ciceroSection{13} What of those luminaries of Marcus Antonius's camp --- do you not set them before your eyes? First, the two colleagues of the Antonii and of Dolabella, Nucula and Lento, the parcellers-out of Italy under that law which the Senate judged to have been carried by violence; of whom the one composed mimes, the other acted a tragedy. What shall I say of Domitius the Apulian, whose goods I lately saw put up for sale? --- so great is the negligence of his agents. And this man recently poured poison into his sister's son: he did not merely give it. But they cannot help living lavishly who count on our goods while they pour out their own. I saw too the auction of Publius Decius, a distinguished man, who, following the examples of his ancestors, devoted himself --- for another man's debt. Yet in that auction not a single buyer was found. A ridiculous fellow, to think he can free himself from his own debt by selling what belongs to others!

\ciceroSection{14} For what shall I say of Trebellius? The Furies of his debtors seem to have had their revenge on him; for the very champion of the new account-books we have seen turned into a new account-book himself. What of Titus Plancus? --- whom that most excellent citizen Aquila drove out of Pollentia, and that with a broken leg: would it had befallen him sooner, that he might not have been able to come back here! I have all but passed over the light and ornament of that army, Titus Annius Cimber, son of Lysidicus --- Lysidicus himself, since he has dissolved every law; unless perhaps it was by right that the Cimbrian killed his own brother, a German. Since Antony has so great an abundance of such men, and of this kind, what crime will he leave undone --- when Dolabella has bound himself by acts of murder so great, with a band and a supply of brigands by no means his equal?

\ciceroSection{15} Wherefore, as I have often dissented from Quintus Fufius against my will, so now I have gladly assented to his motion: from which you ought to judge that I am not in the habit of quarrelling with the man, but with the cause. And so I do not merely assent, I even give thanks to Fufius: for he delivered a severe, weighty motion, worthy of the commonwealth. He judged Dolabella a public enemy; he proposed that his goods be confiscated to the state. And though nothing could have been added to it --- for what more drastic could he have decreed, what more severe? --- yet he declared that, if any of those called on to speak after him should propose a sterner motion, he would go over to it. Who can fail to praise such severity?

\ciceroSection{16} Now, since Dolabella has been judged a public enemy, he must be pursued with war. For he does not keep quiet; he has a legion, he has runaway slaves, he has a criminal band of impious men; he is himself overbold, unrestrained, vowed to a gladiator's manner of death. Wherefore, since Dolabella was decreed an enemy yesterday and war must be waged against him, a commander must be chosen. Two motions have been put, of which I approve neither: the one because I think it always dangerous, except when there is no choice; the other because I judge it ill-suited to these times.

\ciceroSection{17} For an extraordinary command is a popular, windy thing --- least of all befitting our gravity, least of all this order. In the war with Antiochus, great and grievous, when the province of Asia had fallen to Lucius Scipio, and he was thought to have too little spirit in him, too little strength, and the Senate was for transferring the charge to his colleague, Gaius Laelius --- the father of the famous wise man --- Publius Africanus, the elder brother of Lucius Scipio, rose and prayed that disgrace away from his family, and said both that there was the highest virtue and the highest judgment in his brother, and that he himself, at his age and with his record, would not fail him as a legate. When this had been said by him, nothing was changed concerning Scipio's province; nor was an extraordinary command sought for that war any more than in the two greatest Punic wars before it, which were waged and finished by consuls or by dictators, or than for the war with Pyrrhus, or with Philip, or, later, for the Achaean war, or for the Third Punic War; and for that war the Roman people so chose for itself a fit leader, Publius Scipio, that it nonetheless willed him to wage it as consul.

\ciceroSection{18} The war against Aristonicus had to be waged in the consulship of Publius Licinius and Lucius Valerius. The people were asked whom they wished to wage that war. Crassus, the consul and pontifex maximus, set a fine upon his colleague Flaccus, the flamen of Mars, should he depart from the sacred rites: which fine the Roman people remitted; yet they ordered the flamen to obey the pontifex. But not even then did the Roman people commit the war to a private man, although there was Africanus, who the year before had triumphed over the Numantines; and he, though he far surpassed all in the glory and virtue of war, nonetheless carried only two tribes. So the Roman people gave the war to be waged by Crassus the consul rather than by Africanus a private man. As for the commands of Gnaeus Pompeius, that greatest of men and chief of all, it was turbulent tribunes of the plebs who carried them. For the Sertorian war was given by the Senate to a private man, because the consuls refused it, when Lucius Philippus said that he was sending him in place of the consuls --- not as a proconsul.

\ciceroSection{19} What, then, is this electoral assembly, or what canvassing, that Lucius Caesar, that most steadfast and weighty citizen, has brought into the Senate? Upon a most illustrious and blameless man he has decreed a command --- yet to a private man: and in this he has laid the heaviest burden upon us. If I assent, I shall have brought canvassing into the senate-house; if I refuse, I shall seem by my vote, as in an election, to have denied an honour to a man who is my closest friend. But if it is your pleasure to hold elections in the Senate, then let us stand for office, let us canvass --- only let a ballot-tablet be given us, as it is given to the people. Why do you bring it about, Caesar, that either a most excellent man, if he is not assented to, should seem to have suffered a defeat, or each one of us, if passed over --- since we are of equal rank --- should be thought not worthy of the same honour?

\ciceroSection{20} ``But,'' it is objected --- for I hear it whispered --- ``you, by your own motion, gave an extraordinary command to the young Gaius Caesar.'' He, indeed, had given me an extraordinary protection: and when I say ``to me,'' I mean to the Senate and to the Roman people. From a man through whom the commonwealth came to hold a protection so great --- one not even thought of beforehand --- that without it she could not be safe, was I to withhold an extraordinary command? Either the army had to be taken from him, or the command granted. For what reasoning is there, or how can it come about, that an army be held without a command? Therefore what was not snatched away is not to be reckoned as given. You would have snatched the command from Gaius Caesar, senators, had you not granted it. The veteran soldiers, who, following his authority, his command, and his name, had taken up arms for the commonwealth, wished to be commanded by him; the Martian legion and the Fourth had so attached themselves to the authority of the Senate and the dignity of the commonwealth that they demanded Gaius Caesar as their general and leader. To Gaius Caesar the necessity of war gave the command; the Senate gave the rods of office. But to a man at leisure, doing nothing, a private citizen --- I beg you, Lucius Caesar, for it is with a man of the deepest experience that I am dealing --- when did the Senate ever give a command? But on this matter, enough --- lest I should seem to be voting against a man who is my closest friend and has deserved supremely well of me. And yet who can vote against one who not only does not seek the office, but even refuses it?

\ciceroSection{21} But that other motion, senators --- alien to the dignity of the consuls, alien to the gravity of the times --- is the one that the consuls should draw lots for Asia and Syria, for the purpose of pursuing Dolabella. I shall say why it is useless to the commonwealth; but first see how shameful it is for the consuls. When a consul-elect is under siege, when on his deliverance the safety of the commonwealth is set, when pestilent citizens and parricides have broken away from the Roman people, and when we are waging the kind of war in which we contend for our dignity, our liberty, our very life --- so that, should anyone fall into Antony's power, tortures and torments are in store --- and when the deciding of all these matters has been entrusted and committed to consuls the best and bravest, shall mention be made of Asia and Syria, so that we may seem to have furnished either a ground for suspicion or material for ill-will?

\ciceroSection{22} ``But,'' they say, they decree it ``once Brutus has been freed'': for that was all that remained --- that he be left behind, deserted, betrayed. I, for my part, say that mention of the provinces has been made at the most unsuitable time of all. For however intent your mind may be, Gaius Pansa --- as indeed it is --- upon freeing that bravest and most illustrious of all men, yet the nature of things compels you of necessity to turn your mind at some point to the pursuit of Dolabella, and to divert some part of your care and thought to Asia and Syria. But if it could be done, I should wish you to have several minds, that you might bend them all upon Mutina. Since that cannot be, we wish you, with that mind of yours which is most excellent and best, to think of nothing but Brutus.

\ciceroSection{23} This indeed you are doing, and upon it you lean with all your weight, I know; yet no one can either carry on two things at once --- great ones especially --- or even work them out in thought. We ought to spur on and inflame that most excellent zeal of yours, not turn it aside in any part to another care. Add to this men's talk, add their suspicions, add their ill-will: imitate me, whom you have always praised --- me who laid down a province equipped and adorned by the Senate, that I might quench the conflagration of my country, setting all other thought aside. There will be no one --- except me alone, with whom, surely, had you thought it concerned you at all, you would have shared it, given the closeness of our friendship --- who will believe that the province was decreed to you against your will. This rumour, I beg you, by that singular wisdom of yours, suppress; and bring it about that you do not seem to desire the thing you do not care for.

\ciceroSection{24} And on this you must labour the more strenuously, because your colleague, that most illustrious man, cannot fall under the same suspicion. He knows nothing of these things, suspects nothing; he wages war, he stands in the line of battle, he contends for his blood and his breath; he will hear that a province has been decreed to him before he could even suspect that any time had been set aside for that matter. I fear that our armies too --- which betook themselves to the commonwealth not by the compulsion of a levy but by voluntary zeal --- may be slackened in spirit, if they think we have given a thought to anything other than the war at hand. But if provinces seem to be things the consuls should seek --- as they have often been sought by many most illustrious men --- first give us back Brutus, the light and glory of the state; who must be so preserved as that emblem is which, fallen from heaven, is kept under the guardianship of Vesta; while it is safe, we too shall be safe. Then we will lift you, if it can be done, to the very sky upon our shoulders; provinces, surely, the most worthy, we will choose for you; now let us do the thing in hand. And the thing in hand is whether we are to live free or to meet death --- which is surely to be preferred to slavery.

\ciceroSection{25} What more? What if that motion even brings delay to the pursuit of Dolabella? For when will a consul come? Are we to wait until not so much as a trace of the communities and cities of Asia is left? ``But they will send someone of their own number.'' That can win my hearty approval --- I who, a little while ago, refused an extraordinary command to a most illustrious man, a private citizen! ``But they will send a worthy man.'' More worthy than Publius Servilius? Yet such a man the state does not possess. What I myself, then, thought should be given to no one --- not even by the Senate --- shall I approve when it is conferred by the judgment of a single man?

\ciceroSection{26} We have need, senators, of a man unencumbered and ready, and of one who holds a lawful command, who has besides authority, a name, an army, and a spirit proven in liberating the commonwealth. Who, then, is he? Either Marcus Brutus, or Gaius Cassius, or both. I would decree it outright --- as in many decrees, ``the consuls, one or both of them'' --- had we not bound Brutus fast in Greece, and preferred that his aid should incline toward Italy rather than toward Asia: not that we should have something to look back to from the battle-line, but that the battle-line itself should have a reserve even from across the sea. Besides, senators, Marcus Brutus is even now held back by Gaius Antonius, who holds Apollonia, a great and weighty city; he holds, I believe, Byllis, he holds Amantia, he presses upon Epirus, he hangs over Oricum, he has several cohorts, he has cavalry. If Brutus is drawn off from here to another war, we shall surely lose Greece. And provision must also be made for Brundisium and that shore of Italy. Yet I marvel that Antony delays so long; for he is wont to hold out his own wrists for the manacles, and not to bear the dread of a siege any longer. But if Brutus shall have made an end of it, and shall understand that he will profit the commonwealth more by pursuing Dolabella than by remaining in Greece, he will act of his own accord, as he has done thus far also, nor will he, amid so many fires that must be put out at once, wait upon the Senate.

\ciceroSection{27} For both Brutus and Cassius have in many matters already been their own senate. For it is inevitable, amid so great an overturning and disturbance of all things, that one obey the times rather than the precedents. Nor is this the first time that either Brutus or Cassius has judged the safety and liberty of his country the most sacred law and the best of customs. And so, even if nothing were referred to us concerning the pursuit of Dolabella, I should nonetheless count it as good as decreed, since these were such men in virtue, in authority, in nobility --- supreme men, of whom the army of the one was already known to us, that of the other heard of. Did Brutus, then, wait for our decrees, when he knew our wishes? For he did not set out for his own province of Crete: he flew to a province not his own, Macedonia; he reckoned everything his own that you would wish to be yours; he enrolled new legions, took over the old; he drew off to himself the cavalry of Dolabella, and judged him, by his own vote, a public enemy --- though not yet defiled by so great a parricide. For were it not so, by what right would he draw off the cavalry from a consul?

\ciceroSection{28} What of Gaius Cassius, endowed with equal greatness of soul and of judgment --- did he not set out from Italy with this purpose, to keep Dolabella out of Syria? By what law, by what right? By that which Jupiter himself has ordained: that all things which are salutary to the commonwealth be held lawful and just. For law is nothing else than right reason, drawn from the will of the gods, commanding what is honourable and forbidding its opposite. To this law, then, Cassius was obedient when he set out for Syria --- a province not his own, were men governed by written laws, but, those laws once trampled down, his own by the law of nature.

\ciceroSection{29} But that this may be confirmed also by your authority, I propose as follows: Whereas Publius Dolabella, and those who were the agents, the partners, the helpers of his most cruel and most loathsome deed, have been judged by the Senate enemies of the Roman people; and whereas the Senate has resolved that Publius Dolabella be pursued with war, so that he who has polluted all the laws of gods and men by a crime new, unheard of, and inexpiable, and has bound himself by a wicked parricide against his country, may pay to gods and men the penalties he has earned and owes:

\ciceroSection{30} that it is the pleasure of the Senate that Gaius Cassius, as proconsul, hold the province of Syria, as one who has obtained that province by the best right; that he receive the armies from Quintus Marcius Crispus, proconsul, from Lucius Staius Murcus, proconsul, and from Aulus Allienus, legate, and that they hand them over to him; and that with these forces, and any others he may have made ready besides, he pursue Publius Dolabella by land and by sea. For the conducting of that war, that he have the right and power to requisition, from whomever he sees fit, ships, sailors, money, and all else that pertains to the waging of that war, throughout Syria, Asia, Bithynia, and Pontus; and that, into whatever province he shall come for the purpose of waging that war, there the command of Gaius Cassius, proconsul, be greater than that of him who shall then be holding that province, at the time when Gaius

\ciceroSection{31} Cassius, proconsul, shall have come into that province; that King Deiotarus the father and King Deiotarus the son, if --- as in many wars they have time and again aided the dominion of the Roman people --- they shall likewise have aided Gaius Cassius, proconsul, with their forces and their resources, will be doing what is pleasing to the Senate and the Roman people; and likewise, should the other kings, tetrarchs, and princes do the same, that the Senate and the Roman people will not be unmindful of their service. And that Gaius Pansa and Aulus Hirtius, the consuls, one or both, if it seem good to them, once the commonwealth is recovered, refer to this order, at the earliest possible time, the matter of the consular and praetorian provinces. Meanwhile, let the provinces be held by those by whom they are now held, until a successor to each has been appointed by decree of the Senate.

\ciceroSection{32} By this decree of the Senate you will set fire to a man already ablaze, and arm a man already armed --- Cassius; for you can be ignorant neither of his spirit nor of his forces. His spirit is what you see; his forces are what you have heard of --- brave and steadfast men, who, even while Trebonius was yet alive, would not have suffered Dolabella's banditry to penetrate into Syria. Allienus, my friend and close connection, will assuredly, after the death of Trebonius, not even wish to be called the legate of Dolabella. There is the sturdy and victorious army of Quintus Caecilius Bassus --- a private man, indeed, but a brave and distinguished one.

\ciceroSection{33} There is the army of King Deiotarus, both father and son --- a great army, and trained after our own fashion; in the son the highest hope, the highest native talent, the highest virtue. What shall I say of the father? Whose goodwill toward the Roman people is as old as the man himself; who was not only the ally of our commanders in their wars, but even the leader of his own forces. The things which Sulla, which Murena, which Servilius, which Lucullus declared of that man --- how handsomely, how honourably, how weightily, and how often, in the Senate! What of Gnaeus

\ciceroSection{34} Pompeius shall I say --- who judged Deiotarus alone, in all the world, a friend from the heart and truly well-disposed, alone faithful to the Roman people? Marcus Bibulus and I were commanders in neighbouring and bordering provinces: by that same king we were aided, with cavalry and with foot. There followed this most bitter and most calamitous civil war, in which what Deiotarus ought to have done, or what on the whole would have been more right, there is no need to say --- especially since the victory in the war gave its verdict against what Deiotarus thought. If in that war there was error, it was shared by him with the Senate; but if his judgment was right, then not even a defeated cause is to be blamed. To these forces other kings will be added; levies, too, will be added.

\ciceroSection{35} Nor indeed will fleets be wanting: so highly do the Tyrians value Cassius, so great is his name in Syria and in Phoenicia. The commonwealth, senators, has in Gaius Cassius a commander ready against Dolabella --- and not only ready, but skilled and brave. Great were the deeds he did before the arrival of Bibulus, that excellent man, when he routed the noblest leaders of the Parthians and their vast forces, and freed Syria from the monstrous onset of the Parthians. His greatest and most singular praise I pass over; for the proclaiming of it is not yet welcome to all, and we had better preserve it by the testimony of memory rather than of speech.

\ciceroSection{36} I have noticed, senators --- I have even overheard --- that I honour Brutus too much, and that by my motion dominion and supremacy are being given to Cassius. Whom do I honour? Surely those who are themselves the ornaments of the commonwealth. What? Did I not always, in every motion, honour Decimus Brutus? Do you blame me for it, then? Or should I rather honour the Antonii --- the shame and disgrace not of their own families only, but of the Roman name? Should I honour Censorinus, an enemy in war, a dealer in confiscated goods in peace? Should I gather up the rest of the wreckage from the same gang of brigands? For my part, so far am I from honouring those enemies of peace, of concord, of the laws, of the courts, of liberty, that it cannot be brought about that I should not hate them as much as I love the commonwealth.

\ciceroSection{37} ``See,'' he says, ``that you do not offend the veterans'': for this above all I keep overhearing. But I am bound to protect the veterans --- those of them, that is, who are sound of mind; certainly I am not bound to fear them. Those veterans, indeed, who took up arms for the commonwealth and followed Gaius Caesar, the author of their fathers' benefactions, and who today defend the commonwealth at the peril of their lives --- these I am bound not merely to protect, but to enrich with advantages. But those who keep quiet --- like the Seventh legion, like the Eighth --- I think are to be set in great glory and praise. As for the companions of Antony, who, after they had devoured the benefactions of Caesar, now besiege a consul-elect, threaten this city with sword and fire, and have handed themselves over to Saxa and Cafo --- men born for crime and plunder --- is there anyone who thinks they ought to be protected? Men, therefore, are either good, whom we ought even to honour; or quiet, whom we ought to preserve; or impious, against whose madness we have taken up war and lawful arms.

\ciceroSection{38} Whose veterans' feelings, then, are we afraid of offending? Those who long to free Decimus Brutus from his siege? Since the safety of Brutus is dear to them, how can they hate the name of Cassius? Or those who hold aloof from both sets of arms? I do not fear that I shall be harsh to any of those who delight in peace. But upon the third kind --- not veteran soldiers, but most ungovernable enemies --- I long to brand the sharpest pain I can. And yet, senators, how long shall we deliver our motions at the veterans' discretion? What is this disdain of theirs, what arrogance so great, that we should choose even our commanders at their nod?

\ciceroSection{39} For my part --- for I must say what I feel, senators --- I think we have not so much to fear from the veterans as to ask what the raw recruits, the flower of Italy, what the new legions, most ready to liberate their country, what all Italy may feel about your gravity. For nothing flourishes for ever; one age succeeds another: long did the legions of Caesar flourish; now flourish those of Pansa, those of Hirtius, those of the son of Caesar, those of Plancus; they surpass in number, they surpass in youth; nay, they surpass in authority too. For the war they wage is one approved by all nations. And so to these, rewards have been promised; to those, they have been paid in full. Let those enjoy what is theirs; let what we have pledged be paid in full to these. For that, I trust, the immortal gods judge most just.

\ciceroSection{40} Since these things are so, I move, senators, that the motion I have stated be approved by you.
